\noindent Первая часть книги разбирает нейрон как электрическую схему. Это единственный уровень, на котором мы понимаем мозг \emph{количественно}: можем записать уравнения, решить их и сравнить с измерением до десятой доли милливольта. Всё, что лежит выше, --- вычисления, сети, обучение, --- мы будем понимать тем лучше, чем твёрже стоим здесь. Инструментов понадобится немного: закон Ома, законы Кирхгофа, конденсатор, диффузия и больцмановский множитель $e^{-E/kT}$. Всё остальное --- ионные каналы, насосы, медиаторы --- мы введём по дороге как элементы схемы.

План части простой. Сначала --- из каких деталей собран нейрон (глава~\ref{sec:neuron}): изолятор, конденсатор, резисторы, батарейки и зарядное устройство. Потом --- почему на мембране стоит напряжение (глава~\ref{sec:membrane}). Потом --- нервный импульс (глава~\ref{sec:actionpotential}): как нелинейная обратная связь делает из мембраны триггер и как импульс бежит по аксону. Потом --- синапс (глава~\ref{sec:synapse}): как импульс передаётся от выхода одного нейрона на вход другого. И в конце --- дендрит и аксон как линии передачи (глава~\ref{sec:cable}).
